\documentclass[11pt]{article}
\usepackage[margin=1in]{geometry}
\usepackage{hyperref, booktabs, array}
\usepackage[shortlabels]{enumitem}
\usepackage{fancyhdr}
\usepackage{microtype}

% -------------------------------------------------------
\newcommand{\COURSE}{ECON [XXX]}
\newcommand{\COURSETITLE}{[Course Title]}
\newcommand{\TERM}{[Term Year]}
\newcommand{\INSTRUCTOR}{[Instructor Name]}
\newcommand{\EMAIL}{[email@university.edu]}
\newcommand{\OFFICEHOURS}{[Day, Time, Location]}
\newcommand{\LECTURES}{[Day/Time, Building/Room]}

\pagestyle{fancy}
\fancyhf{}
\lhead{\COURSE{} --- \TERM}
\rhead{Syllabus}
\cfoot{\thepage}

% ============================================================
\begin{document}

\begin{center}
  {\LARGE\bfseries \COURSE: \COURSETITLE}\\[0.4em]
  {\large \TERM}
\end{center}
\medskip\hrule\medskip

% -------------------------------------------------------
\section*{Course Information}

\begin{tabular}{@{}ll}
  \textbf{Instructor:}    & \INSTRUCTOR \\
  \textbf{Email:}         & \href{mailto:\EMAIL}{\EMAIL} \\
  \textbf{Office Hours:}  & \OFFICEHOURS \\
  \textbf{Lectures:}      & \LECTURES \\
\end{tabular}

% -------------------------------------------------------
\section*{Course Description}

[Brief course description: what students will learn and why it matters.]

\textbf{Prerequisites:} [List prerequisites.]

% -------------------------------------------------------
\section*{Learning Outcomes}

By the end of this course, students will be able to:
\begin{enumerate}
  \item [Outcome 1.]
  \item [Outcome 2.]
  \item [Outcome 3.]
\end{enumerate}

% -------------------------------------------------------
\section*{Required Materials}

\begin{itemize}
  \item \textit{[Textbook Title]}, [Author(s)], [Edition].
  \item Additional readings posted on course website.
\end{itemize}

% -------------------------------------------------------
\section*{Assessment}

\begin{center}
\renewcommand{\arraystretch}{1.3}
\begin{tabular}{lrr}
  \toprule
  \textbf{Component}   & \textbf{Weight} & \textbf{Due} \\
  \midrule
  Problem Sets (X)     & 20\%            & Biweekly \\
  Quizzes              & 10\%            & In class \\
  Midterm 1            & 20\%            & [Date] \\
  Midterm 2            & 20\%            & [Date] \\
  Final Exam           & 30\%            & [Date] \\
  \midrule
  \textbf{Total}       & \textbf{100\%}  & \\
  \bottomrule
\end{tabular}
\end{center}

% -------------------------------------------------------
\section*{Grading Scale}

\begin{tabular}{llll}
  A: 90--100\% & B: 80--89\% & C: 70--79\% & D: 60--69\%
\end{tabular}

% -------------------------------------------------------
\section*{Course Schedule}

\renewcommand{\arraystretch}{1.3}
\begin{tabular}{>{\bfseries}cl}
  \toprule
  Week & Topic \\
  \midrule
  1  & Introduction and Overview \\
  2  & [Topic] \\
  3  & [Topic] \\
  4  & [Topic] \\
  5  & [Topic] \\
  6  & [Topic] --- \textbf{Midterm 1} \\
  7  & [Topic] \\
  8  & [Topic] \\
  9  & [Topic] \\
  10 & [Topic] \\
  11 & [Topic] \\
  12 & [Topic] --- \textbf{Midterm 2} \\
  13 & [Topic] \\
  14 & Review \\
     & \textbf{Final Exam: [Date, Time, Room]} \\
  \bottomrule
\end{tabular}

% -------------------------------------------------------
\section*{Policies}

\textbf{Late Work:} [Policy on late submissions.]

\textbf{Collaboration:} Problem sets may be worked on in groups of up to [N]. Exams are individual.

\textbf{Academic Integrity:} All submitted work must be your own. Violations will be handled according to university policy.

\textbf{Accessibility:} Students requiring accommodations should contact the accessibility office and provide documentation within the first two weeks.

\end{document}
